% 第五章内容结构与建模路线图（论文图 2 的源文件）
% 依据：paper/sections/05_model_and_solution.tex 的标题层级
%       5.1 / 5.2.1-5.2.3 / 5.3.1-5.3.6
% 编译：xelatex -interaction=nonstopmode -output-directory=paper/figures code/figures/fig13_chapter5_roadmap.tex
% 产物：paper/figures/fig13_chapter5_roadmap.pdf
% 约定：框内只用简短总结语，不加任何公式；无底色、无分节线，纯黑白线框。
%       三处分支拆成"物性／时程／区域"三框后再合并；定解条件的三个条件嵌套在
%       "定解条件"容器框内。
\documentclass[border=8pt,12pt]{standalone}
\usepackage{xeCJK}
\usepackage{tikz}
\usetikzlibrary{shapes.geometric,arrows.meta}
\setCJKmainfont{SimSun}
\setCJKfamilyfont{zhhei}{SimHei}
\setmainfont{Times New Roman}
\pagestyle{empty}

\tikzset{
  base/.style = {align=center, font=\small\CJKfamily{zhhei}, draw=black, fill=white,
                 line width=0.9pt, rounded corners=3pt, inner sep=3pt},
  bx/.style   = {base, rectangle, text width=2.9cm, minimum height=1.6cm},
  bxs/.style  = {base, rectangle, rounded corners=2pt, line width=0.7pt,
                 text width=2.5cm, minimum height=0.9cm},
  bxi/.style  = {base, rectangle, rounded corners=2pt, line width=0.7pt,
                 minimum width=2.2cm, minimum height=0.85cm, inner sep=2pt},
  dec/.style  = {base, diamond, rounded corners=1.5pt, aspect=1.95,
                 minimum height=1.9cm, inner sep=0pt, line width=1.0pt},
  term/.style = {base, rounded corners=0.8cm,
                 minimum height=1.6cm},
  ar/.style   = {-{Latex[length=5pt,width=4.5pt]}, line width=0.9pt, rounded corners=6pt},
  ln/.style   = {line width=0.9pt, rounded corners=6pt},
  lab/.style  = {font=\scriptsize\CJKfamily{zhhei}, inner sep=2pt, fill=white},
  ctitle/.style = {font=\small\CJKfamily{zhhei}, anchor=center, inner sep=0pt},
}

\begin{document}
\begin{tikzpicture}

% ============ 建模准备：统一模型 → 三处分支 → 合并 ============
\node[term, text width=3.2cm] (t1) at (2.50,-1.85) {统一径向轴对称模型};
\node[bxs] (p1) at (7.00,-0.80) {物性};
\node[bxs] (p2) at (7.00,-1.85) {时程};
\node[bxs] (p3) at (7.00,-2.90) {区域};
\node[bx]  (b3) at (11.90,-1.85) {守恒分析得\\两方程};

\coordinate (sp) at (4.90,-1.85);
\coordinate (mg) at (9.40,-1.85);
\draw[ln] (t1.east) -- (sp);
\draw[ar] (sp) |- (p1.west);
\draw[ar] (sp) -- (p2.west);
\draw[ar] (sp) |- (p3.west);
\draw[ln] (p1.east) -| (mg);
\draw[ln] (p2.east) -- (mg);
\draw[ln] (p3.east) -| (mg);
\draw[ar] (mg) -- (b3.west);

% ============ 定解条件（容器框，内嵌三个条件框） ============
\node[base, rectangle, rounded corners=4pt, minimum width=8.2cm, minimum height=2.1cm,
      inner sep=0pt] (cond) at (4.80,-5.35) {};
\node[ctitle] at (4.80,-4.67) {定解条件};
\node[bxi] (c1) at (2.10,-5.80) {边界条件};
\node[bxi] (c2) at (4.70,-5.80) {初值};
\node[bxi] (c3) at (7.30,-5.80) {环境驱动};
\node[bx]  (b5) at (11.90,-5.35) {量级判定\\Fourier／Biot};

\draw[ar] (b3.south) -- (11.90,-3.70) -- (4.80,-3.70) -- (cond.north);
\draw[ar] (cond.east) -- (b5.west);

% ============ 求解主线 ============
\node[bx]  (b6) at (2.30,-8.20)  {顶点有限\\体积离散};
\node[bx]  (b7) at (6.30,-8.20)  {端点处理与\\时间推进};
\node[dec] (b8) at (10.20,-8.20) {含移动\\边界？};
\node[bx]  (b9) at (14.20,-8.20) {坐标变换与\\表观对流项};

\node[bx]  (b10) at (14.20,-10.80) {守恒核验与\\三路验证};
\node[term, text width=3.0cm] (b11) at (7.50,-10.80) {输出温湿场};

\draw[ar] (b5.south) -- (11.90,-6.85) -- (2.30,-6.85) -- (b6.north);
\draw[ar] (b6) -- (b7);
\draw[ar] (b7) -- (b8);
\draw[ar] (b8.east) -- node[lab, above]{是} (b9.west);
\draw[ar] (b9.south) -- (b10.north);

% 不含移动边界时直接进入核验
\draw[ar] (b8.south) -- (10.20,-9.55) -- node[lab, above]{否} (12.90,-9.55)
                      -- (12.90,-10.00);

% 回环：三路验证不一致则加密网格
\draw[ar] (b10.south) -- (14.20,-12.40) -- node[lab, above]{不一致·加密} (7.00,-12.40)
                        -- (0.25,-12.40) -- (0.25,-8.20) -- (b6.west);

\draw[ar] (b10) -- (b11);

\end{tikzpicture}
\end{document}
